\documentclass[10pt,a4paper]{amsart}
\usepackage[margin=19mm]{geometry}
\usepackage{amsmath,amssymb,mathtools}
\usepackage[scr=rsfs,cal=euler]{mathalfa}
\usepackage{xfrac}
\usepackage[unicode,hidelinks,bookmarks=false]{hyperref}
\newcommand{\ie}{i.e.\ }
\newcommand{\cf}{cf.\ }
\newcommand{\resp}{respectively\ }
\newcommand{\Subsection}{\subsection}
\providecommand{\nobreakdash}{\nobreak\hbox{-}}
\setlength{\emergencystretch}{2em}
\multlinegap0pt
\usepackage{amssymb}
\usepackage{tikz}
\newtheorem{prop}{Proposition}
\title{Ignorance with(out) Grasping}
\author{Ekaterina Kubyshkina, Mattia Petrolo}
\date{Source: arXiv:2603.09569v1 (CC BY 4.0)}
\begin{document}
\maketitle

\noindent\emph{Source and licence.} Ignorance with(out) Grasping. Version arXiv:2603.09569v1 (CC BY 4.0), distributed by arXiv under the Creative Commons Attribution 4.0 International licence, http://creativecommons.org/licenses/by/4.0/, as recorded in the arXiv licence metadata for this exact version and shown on the abstract page https://arxiv.org/abs/2603.09569v1. Full licence notice: \texttt{licenses/arxiv-2603.09569-CC-BY-4.0.txt}.\par\smallskip

\noindent\textit{Editorial scope.} Counted unit: the article's prop genA6IW (source0.tex lines 843--854). The article's complete original proof of it is delivered with it (source0.tex lines 856--896, one \texttt{proof} environment of 41 lines). No mathematics is written, repaired or re-derived: the statement and the proof are the article's own lines, in the article's own notation, and no retained line is substituted or re-flowed.  No further source-local context is needed: every cross-reference of every retained line resolves inside the two delivered spans.  Nothing was omitted from any retained span, so the package carries no omission record.  No retained line contains a citation, so the wrapper carries no bibliography and no bibliography entry.  No retained line contains a figure environment, an image-inclusion command or drawing code, so no image is delivered and the package declares no asset dependency.  Editorial formatting only, in addition: an AMS \texttt{amsart} preamble that restates the article's own declarations and macros used by the retained lines, namely the environments \texttt{prop} on separate counters, together with the standard packages those lines need.  The retained environments print in this document as Proposition 1 in order of appearance, whereas the article prints the same items with the numbering of its own full document.  The complete native source of the article as distributed in the arXiv e-print for this exact version is bundled unchanged as \texttt{sources/196071/source0.tex}.
\begin{prop}
\label{genA6IW}

For all $k \geq 1:$

\begin{center}

$\vdash (\neg I^{w} (\bigwedge^{k}_{j=1} \phi_{j} \rightarrow \neg \psi) \wedge \bigwedge^{k}_{j=1} \neg I^{w} \phi_{j} \wedge \bigwedge^{k}_{j=1} \neg I^{w}(\psi \rightarrow \phi_{j})) \rightarrow \neg I^{w} \psi$

\end{center}

\end{prop}

\begin{proof}

\

\begin{enumerate}

\item $(\neg I^{w} (\bigwedge^{k}_{j=1} \phi_{j} \rightarrow \neg \psi) \wedge \neg I^{w}\bigwedge^{k}_{j=1} \phi_{j}  \wedge \neg I^{w}(\psi \rightarrow \bigwedge^{k}_{j=1} \phi_{j} )) \rightarrow \neg I^{w} \psi$ (by (A6$_{I^{w}}$));

\item $\bigwedge^{k}_{j=1} \neg I^{w} \phi_{j} \rightarrow \neg I^{w} \bigwedge^{k}_{j=1} \phi_{j}$ (from (A3$_{I^{w}}$) applied $k - 1$ times);

\item $\bigwedge^{k}_{j=1} \neg I^{w}(\psi \rightarrow \phi_{j}) \rightarrow \neg I^{w} \bigwedge^{k}_{j=1} (\psi \rightarrow \phi_{j})$ (from (A3$_{I^{w}}$) applied $k-1$ times);

\item $\bigwedge^{k}_{j=1} (\psi \rightarrow \phi_{j}) \leftrightarrow (\psi \rightarrow \bigwedge^{k}_{j=1} \phi_{j})$ (TAUT);

\item $(\neg I^{w} \overline{\bigwedge^{k}_{j=1} (\psi \rightarrow \phi_{j})} \wedge \neg I^{w} \overline{(\psi \rightarrow \bigwedge^{k}_{j=1} \phi_{j})}) \rightarrow (\neg I^{w} \bigwedge^{k}_{j=1} (\psi \rightarrow \phi_{j}) \leftrightarrow \neg I^{w} (\psi \rightarrow \bigwedge^{k}_{j=1} \phi_{j}))$ (from 4, by (RE$_{I^{w}}$));

\item $\neg I^{w} \overline{\bigwedge^{k}_{j=1} (\psi \rightarrow \phi_{j})} \leftrightarrow \neg I^{w} \overline{(\psi \rightarrow \bigwedge^{k}_{j=1} \phi_{j})}$ (by definition of $\bar{\chi_{1}}$ and $\bar{\chi_{2}}$ for any $\chi_{1}$ and $\chi_{2}$ in which occur only the same propositional variables);

\item $\neg I^{w} \bigwedge^{k}_{j=1} (\psi \rightarrow \phi_{j}) \rightarrow \neg I^{w} \overline{\bigwedge^{k}_{j=1} (\psi \rightarrow \phi_{j})}$ (from (A5$_{I^{w}}$) by contraposition);

\item $\neg I^{w} \bigwedge^{k}_{j=1} (\psi \rightarrow \phi_{j}) \rightarrow (\neg I^{w} \overline{\bigwedge^{k}_{j=1} (\psi \rightarrow \phi_{j})} \wedge \neg I^{w} \overline{(\psi \rightarrow \bigwedge^{k}_{j=1} \phi_{j})})$ (from 6, 7, by CPL);
 

\item $\neg I^{w} \bigwedge^{k}_{j=1} (\psi \rightarrow \phi_{j}) \rightarrow (\neg I^{w} \bigwedge^{k}_{j=1} (\psi \rightarrow \phi_{j}) \leftrightarrow \neg I^{w} (\psi \rightarrow \bigwedge^{k}_{j=1} \phi_{j}))$ (from 5, 8, by transitivity);

\item $\neg I^{w} \bigwedge^{k}_{j=1} (\psi \rightarrow \phi_{j}) \rightarrow \neg I^{w} (\psi \rightarrow \bigwedge^{k}_{j=1} \phi_{j})$ (from 9, by CPL);

\item $\bigwedge^{k}_{j=1} \neg I^{w}(\psi \rightarrow \phi_{j}) \rightarrow \neg I^{w} (\psi \rightarrow \bigwedge^{k}_{j=1} \phi_{j})$ (from 3, 10, by transitivity);

\item $\neg I^{w} (\bigwedge^{k}_{j=1} \phi_{j} \rightarrow \neg \psi) \rightarrow \neg I^{w} (\bigwedge^{k}_{j=1} \phi_{j} \rightarrow \neg \psi)$ (TAUT);

\item $(\neg I^{w} (\bigwedge^{k}_{j=1} \phi_{j} \rightarrow \neg \psi) \wedge \bigwedge^{k}_{j=1} \neg I^{w} \phi_{j} \wedge \bigwedge^{k}_{j=1} \neg I^{w}(\psi \rightarrow \phi_{j})) \rightarrow (\neg I^{w} (\bigwedge^{k}_{j=1} \phi_{j} \rightarrow \neg \psi) \wedge \neg I^{w}\bigwedge^{k}_{j=1} \phi_{j}  \wedge \neg I^{w}(\psi \rightarrow \bigwedge^{k}_{j=1} \phi_{j} ))$ (from 2, 11, 12, by CPL);

\item $(\neg I^{w} (\bigwedge^{k}_{j=1} \phi_{j} \rightarrow \neg \psi) \wedge \bigwedge^{k}_{j=1} \neg I^{w} \phi_{j} \wedge \bigwedge^{k}_{j=1} \neg I^{w}(\psi \rightarrow \phi_{j})) \rightarrow \neg I^{w} \psi$ (from 1, 13, by transitivity).


\end{enumerate}



\end{proof}

\end{document}
